% Seifert-van Kampen pushout square in the category of groupoids.
\begin{tikzpicture}[
  every node/.style={font=\small},
  arrow/.style={-{Stealth[length=3mm]}, thick},
]
  \node (UV) at (0, 2) {$\pi_1(U \cap V,\; A \cap U \cap V)$};
  \node (U)  at (6, 2) {$\pi_1(U,\; A \cap U)$};
  \node (V)  at (0, 0) {$\pi_1(V,\; A \cap V)$};
  \node (X)  at (6, 0) {$\pi_1(X,\; A)$};

  \draw[arrow] (UV) -- node[above] {$i_{U*}$} (U);
  \draw[arrow] (UV) -- node[left]  {$i_{V*}$} (V);
  \draw[arrow] (U)  -- node[right] {$j_{U*}$} (X);
  \draw[arrow] (V)  -- node[below] {$j_{V*}$} (X);

  \node[font=\small\itshape] at (3, 1) {pushout};

  \begin{scope}[on background layer]
    \fill[orange!5, rounded corners=12pt] (-2.7, -0.9) rectangle (8.7, 3.0);
  \end{scope}

  \node at (3, -1.6)
    {\itshape The fundamental groupoid sees opens, intersections, and the union as one pushout.};
\end{tikzpicture}
